\chapter{Discussion} \label{Chap5}\label{chap:discussion}

\section{Interpretation of the research scope}

This dissertation examines how overlapping satellite-conditioned predictions
can form a continuous coloured regional surface and then become editable urban
assets.  It contributes at two interfaces: density-domain fusion and colour
recovery around the released Sat3DGen model, followed by geospatial association
and publication through OSM, DSM, and ChordAtlas.  This addresses the gap
between plausible window-level generation and spatially identified, reusable
urban content.

\section{Research question: extent and reliability boundary}

The study establishes a working large-image reconstruction and
coordinate-explicit integration at building-asset scale.  Overlapping density
windows are fused into one surface, per-vertex appearance is recovered in a
second pass, and the resulting mesh can be transformed, associated with OSM
footprints, processed against DSM context, and converted for ChordAtlas.
Stable selection identities, manifests, cache checks, and validation at the
Python--Java boundary make these transformations inspectable rather than
implicit.

Reliable regional operation depends on four properties highlighted by the
results.  Overlapping predictions need continuous placement and smooth
weighting; spatial predicates must prevent unrelated or
vertically separated surfaces from being combined; geometry processing must
preserve topology and apply external height information selectively; and
publication must be acknowledged by the consuming application.  These
requirements connect large-area reconstruction quality to asset-level system
reliability.

\section{Contribution of the ChordAtlas integration}

The contribution is defined by the new data and execution contracts introduced
at the ChordAtlas boundary.  The integration extends the existing GIS workflow
in the following ways:

\begin{itemize}
  \item a coordinate contract that preserves geographic selection, a local
  origin, and building identity across mesh conversion;
  \item a selectable large-image source mode that binds the mosaic plan,
  density fusion, colour recovery, coordinate conversion, and footprint crop;
  \item manifest-driven planning and cache validation that bind a request to
  its required tiles and processing stages;
  \item OSM- and DSM-aware processing that associates generated surfaces with
  footprints and publishes separately editable building assets;
  \item selection-scoped roof references and panoramas, exercised for the
  independent-tile selections, that connect appearance inputs to the same
  footprint identity and local coordinate frame;
  \item staged Python--Java publication with explicit failure states,
  background execution, and a dedicated per-vertex-colour renderer; and
  \item executable checks for coordinate conversion, mesh prerequisites,
  selection manifests, cache state, and publication contracts.
\end{itemize}

Together, these changes turn a fused or independently assembled surface into
geographically identified ChordAtlas assets whose processing history and
failure state can be examined.  Sat3DGen provides learned geometry and
appearance, while ChordAtlas provides asset structure and procedural editing.
The bridge preserves the coloured OBJ for display and builds MiniMesh as a
separate workspace representation; successful colour transfer should not be
conflated with colour being embedded in the MiniMesh tiles themselves.

For the independent-tile selections, the roof-reference and panorama results
are important because they associate
appearance observations with the same building identity used by the geometry
pipeline.  They provide registered inputs for ChordAtlas style generation.
The separately retained appearance exports establish the availability of
generated appearance artefacts, while reference registration, material
completeness, and texture fidelity
remain distinct evaluation questions.

\section{Relationship among the evaluated systems}

The project contributes at two interfaces to Sat3DGen.  First, it extends the
released large-image wrapper while retaining the pretrained DINOv3/tri-plane
representation and one global density-domain surface extraction
\cite{qian2026sat3dgen,qianSat3DGenCode2026}.  Fractional placement and smooth
overlap weighting address quantisation and hard contributor changes, while a
second pass restores the colour capability available in single-image
inference.  The matched hard-box experiment isolates the effect of this fusion
rule on boundary continuity and resource use.

Secondly, the mesh-space path introduces geographic placement, OSM/DSM
constraints, building extraction, and ChordAtlas publication.  This path
remains relevant where independently extracted or cached meshes are the
available interface.  The two experiments therefore compare representation
stages rather than treating mesh stitching as the only regional strategy.
Independent reference geometry and street-level imagery are still required to
evaluate metric and perceptual accuracy.

The retained London examples clarify that the systems address different
subproblems.  CityEngine generates controllable polygonal content from GIS
shapes, attributes, assets, and CGA rules and supports scripted procedural
runtimes \cite{esriCityEngineIntro2026,esriCityEngineGIS2026,
esriCityEngineSDK2026}.  Its retained export is a compact UV-mapped envelope
whose openings are represented primarily in repeated facade atlases.  Sat3DGen
instead conditions dense geometry and vertex appearance on overhead imagery.
The retained ChordAtlas exports organise surfaces into UV material groups and
appearance maps.  The bridge connects these representations through stable
identity, coordinates, and editable asset boundaries.

Table~\ref{tab:paradigm-comparison} summarises the distinct evidence and role of
each system in this study.

\begin{table}[H]
\centering
\caption{Evaluated systems, output representations, and their roles in the
study.}
\label{tab:paradigm-comparison}
\footnotesize
\begin{tabularx}{\textwidth}{>{\raggedright\arraybackslash}p{0.16\textwidth}
>{\raggedright\arraybackslash}X>{\raggedright\arraybackslash}X
>{\raggedright\arraybackslash}X}
\toprule
System & Input and control & Output representation & Evidence and role\\
\midrule
CityEngine / CGA & Platform support for GIS shapes, CGA rules, attributes, and
assets; one initial shape in the retained case & Compact UV-mapped polygon
export & Illustrative export for compactness and material organisation\\
Sat3DGen large-image & Rectangular mosaic, overlap schedule, and fusion weights
& One density-fused surface with predicted vertex RGB & Directly matched OSM
case; evaluated for continuity, colour support, and topology\\
Sat3DGen independent tiles & Nine image-conditioned meshes plus OSM/DSM
operations & Stitched per-footprint assets with stored vertex RGB & Directly
matched OSM case; evaluated for mesh-space assembly behaviour\\
ChordAtlas procedural appearance & Structured procedural surfaces and style
controls & Semantic UV groups with diffuse, normal, and specular maps &
Retained procedural appearance representation; not reconstruction truth or topology
preservation\\
Hybrid bridge & Geographic selection, manifests, OSM/DSM context, and
publication controls & Identified coloured OBJ, MiniMesh workspace, and
registered appearance inputs & Integration contribution across learned and
procedural representations\\
\bottomrule
\end{tabularx}
\end{table}

This is a partially matched representation case rather than a common reconstruction
benchmark.  The two Sat3DGen routes share bounds and footprint identities,
whereas the CityEngine export lacks the common OSM identifier and ChordAtlas
changes the asset representation.  Topology and colour-variation statistics
are therefore compared only between the two Sat3DGen paths.  Across systems,
the evidence supports conclusions about compactness, material organisation,
available control mechanisms, and workflow structure.

\section{Architectural implications of the manifest-driven design}

The main value of the manifests is that they turn hidden assumptions into
explicit system state.  A request can record the intended geographic extent,
required inputs, permitted tiles, completed stages, and reason for failure.
This enables incomplete acquisition or processing to stop visibly instead of
leaving an output directory that appears complete.  Exact tile selection is
especially important in a shared cache, where recursively consuming every
available mesh can combine unrelated locations while still producing
syntactically valid output.

However, a manifest is authoritative only when it binds outputs to their full
provenance.  Imagery, coordinate reference system, model checkpoint,
inference options, processing configuration, code revision, and output hashes
all influence the meaning of a cached result.  The current design captures
several of these fields but does not yet enforce all of them consistently.
Stable geometric identity is therefore useful for request reuse, but it must be
paired with configuration and source identity before cache reuse can imply
reproducibility.

The cross-language boundary introduces a related transaction problem.  Python
can finish publication before Java has converted and loaded the corresponding
asset. A stronger design would represent processing, conversion, and
visible loading as separate acknowledged stages and would commit the final
state only after the user confirms success.

\section{Implications of the reliability findings}

The large-image comparison shows that nominal overlap alone does not guarantee
a continuous surface.  With hard cropping and equal weights, the set of
contributing windows changes abruptly at regular intervals, turning local
prediction disagreement into visible height steps.  Fractional placement and
raised-cosine weighting substantially reduce those steps without increasing
the measured GPU peak, at the cost of modest additional runtime.  Feathering
smooths disagreement rather than resolving its cause; broad height bias shared
by a window can therefore remain as a gradual surface variation.

The data audit shows that structural completeness and semantic validity are
different properties.  Files may be present while containing duplicates,
cross-partition leakage, inconsistent coordinates, or candidates outside their
intended spatial support.  Dataset preparation should therefore validate
identity, containment, and geographic separation in addition to file counts
and schema shape.  The London workflow is useful as prepared project data, but
these findings prevent treating it as a validated training benchmark.

The mesh experiments reveal the same distinction between execution and
correctness.  A deterministic merge can still select non-contiguous inputs;
a configured vertical threshold can remain ineffective if it is not applied
by the matching predicate; and a DSM correction can produce a complete file
while displacing surfaces that should not be treated as terrain.  The relevant
engineering principle is that every geometric assumption should have both an
executable invariant and a regression case that can falsify it.

Building extraction also creates a trade-off between modularity and geometric
quality.  Footprint association produces independently addressable assets,
which is valuable for editing and procedural processing, but cropping and
capping do not by themselves guarantee closed, manifold, or semantically
separated buildings.  Topology and surface accuracy therefore require explicit
post-processing and evaluation rather than being inferred from successful OBJ
publication.

These observations lead to four general requirements for reliable integration:

\begin{itemize}
  \item validate spatial and semantic relationships, not only the presence and
  shape of files;
  \item make tile membership and provenance authoritative at every stage;
  \item test every geometric assumption and specify how conflicting attributes
  are combined.
\end{itemize}

\section{Contribution assessment}

The contribution is a hybrid reconstruction and integration method rather than
a replacement for either learned reconstruction or procedural modelling.  It
extends the released regional wrapper with continuous density and colour
fusion, then connects the resulting surface to a shared spatial frame,
building identity, external geographic context, modular publication, and
executable data-transfer contracts.  The London data analysis and controlled
geometry tests expose assumptions that would otherwise remain hidden at these
interfaces.

The project builds on Sat3DGen, ChordAtlas, ProceX, BigSUR, and FrankenGAN
\cite{qian2026sat3dgen,kellyChordAtlas2024,kelly2011procex,
kelly2017bigsur,kelly2018frankengan}.  The project-specific contribution is the
London data workflow, large-image density and colour method, geospatial
mesh-processing method, ChordAtlas bridge, and evaluation developed around
these foundations.

\section{Threats to validity}

\subsection{Construct validity}

Most reported measures concern structure and process: coordinates, manifests,
vertices, faces, boundary discontinuity, topology, cache state, and publication
state.  These measures show whether transformations and contracts operate as
specified, but they are not substitutes for geometric or perceptual accuracy.
In particular, smoother seams, complete colour coverage, footprint overlap,
and successful publication do not by themselves establish fa\c{c}ade fidelity
or visual realism.

\subsection{Internal validity and provenance}

The large-image records bind the input plan, script hashes, processing
parameters, and output hashes, but identify the model by repository name rather
than one immutable revision shared explicitly by the density and colour
stages.  The current deployment also resolves those scripts through configured
local paths.  Packaging the scripts with the bridge and binding both stages to
one model revision would make the experiment portable.  At the consumer
boundary, Python requires RGB for the large-image contract, whereas Java can
still fall back from colour parsing; the same requirement should be enforced
at both sides of the transaction.

\subsection{External and statistical validity}

The evaluation is concentrated in one London workflow.  The fusion comparison
uses a matched input, but city, morphology, acquisition, and checkpoint were
not varied.  External services also introduce network, credential, and
licensing dependencies.  Broader claims require a versioned multi-area
evaluation with independent geometry and appearance references.

\section{Priorities for further work}

Three experiments follow directly from the findings.  First, the matched
fusion study should be repeated across multiple urban morphologies and model
revisions, with independent surface and street-level colour references.  A
release-default hard-box run, the fractional-feather method, and relevant
weighting ablations would separate boundary continuity from metric and
perceptual accuracy.

Secondly, the global density path and the independent-mesh path should be run
over the same area and coordinate contract.  Revised vertical stitching and
DSM correction can then be assessed stage by stage for topology, height error,
and surface preservation rather than compared through unrelated cached inputs.

Thirdly, the model revision, scripts, RGB requirement, MiniMesh conversion,
and visible loading acknowledgement should become one portable publication
contract.  The resulting assets can then support a matched CityEngine/CGA
authoring study that extends the present representation example with shared
identity, inputs, and edit tasks, followed by a multi-user Ubiq study of
interactive selection and editing.
